\chapter{Data Structure Dichotomies (D9--D14)}
\label{ch:ds-dichotomies}

\section{D9: Stack $\leftrightarrow$ Recursion}

\textbf{Paradigm A} --- recursive tree traversal:
\begin{lstlisting}
def f(tree):
    if tree is None: return 0
    return 1 + f(tree.left) + f(tree.right)
\end{lstlisting}

\textbf{Paradigm B} --- explicit stack:
\begin{lstlisting}
def g(tree):
    count = 0
    stack = [tree]
    while stack:
        node = stack.pop()
        if node is None: continue
        count += 1
        stack.append(node.left)
        stack.append(node.right)
    return count
\end{lstlisting}

\textbf{Path construction.}
Both compute the tree catamorphism with $\mathsf{base} = 0$ and
$\mathsf{combine}(l, r, \_) = 1 + l + r$.  The recursive version
uses the call stack; the iterative version uses an explicit
$\mathsf{Pair}$-chain as stack.

The compiler produces:
\begin{itemize}
  \item A: $\mathsf{tree\_rec}(\mathsf{tree\_rec\_f}, p_0)$ ---
    a tree RecFunction.
  \item B: $\mathsf{wh\_count}(0, [p_0])$ --- a while-loop RecFunction
    over stack state.
\end{itemize}

\textbf{Automatic inference.}
The path is inferred when:
\begin{enumerate}
  \item The recursive version's RecFunction matches the tree
    catamorphism pattern (base case on None/leaf, recursive calls
    on children).
  \item The iterative version's while-loop processes a stack by
    popping and pushing children --- the compiler detects the
    append/pop pattern on the stack variable.
  \item The $\mathsf{combine}$ function is commutative ($+$ here),
    so traversal order doesn't matter (BFS/DFS/any order).
\end{enumerate}

Certificate: $\mathsf{D9}(\mathsf{tree\_cata}, +, 0) \cdot
\mathsf{add\_comm}$.

\section{D10: \texttt{heapq} $\leftrightarrow$ \texttt{bisect}}

\textbf{Paradigm A}: Kahn's toposort using \texttt{heapq.heappush}/\texttt{heappop}.

\textbf{Paradigm B}: Kahn's toposort using \texttt{bisect.insort}/\texttt{pop(0)}.

\textbf{Path construction.}
Both maintain a sorted frontier of zero-indegree nodes.  Both extract
the minimum element at each step.  The path uses the \emph{ADT
abstraction}: any implementation of the priority queue interface
$(\mathsf{insert}, \mathsf{extract\_min})$ that maintains the
sorted invariant produces the same sequence of extracted elements.

The compiler maps both to the same shared function symbols:
$\mathsf{py\_meth\_heappush}$ and $\mathsf{py\_call\_bisect\_insort}$
both map to $\mathsf{py\_pq\_insert}$; $\mathsf{py\_meth\_heappop}$
and \texttt{.pop(0)} both map to $\mathsf{py\_pq\_extract\_min}$.

\textbf{Automatic inference.}
The compiler maintains a table of \emph{ADT equivalence classes}:
operations that implement the same abstract interface.  When both
programs use operations from the same equivalence class on the
same data, they get the same shared symbols.

\section{D11: \texttt{OrderedDict} $\leftrightarrow$ LinkedList LRU}

\textbf{Paradigm A}: LRU cache using \texttt{OrderedDict.move\_to\_end}.

\textbf{Paradigm B}: LRU cache using a doubly-linked list + dict.

\textbf{Path.}
Both implement the LRU interface: $\mathsf{get}(k)$ (returns value,
moves to front), $\mathsf{put}(k, v)$ (inserts or updates, evicts
LRU if full).  Both shared-symbol to $\mathsf{py\_lru\_get}$,
$\mathsf{py\_lru\_put}$.  Same interface $\Rightarrow$ same behavior
on the same sequence of operations (by ADT abstraction).

\section{D12: \texttt{defaultdict} $\leftrightarrow$ \texttt{setdefault}}

Both implement ``get existing value or create default'':
\begin{lstlisting}
# A: d = defaultdict(list); d[k].append(v)
# B: d.setdefault(k, []).append(v)
\end{lstlisting}

Both compile to: $\mathsf{py\_mut\_append}(\mathsf{py\_meth\_setdefault}(d, k,
\mathsf{NoneObj}), v)$ (same shared symbols).  Path: $\mathsf{refl}$.

\section{D13: \texttt{Counter} $\leftrightarrow$ Manual Counting}
\label{sec:d13-counter}

\texttt{Counter(xs)} is syntactic sugar for
$\fold(\lambda d,x.\, \mathsf{py\_meth\_\_\_setitem\_\_}(d, x,
\mathsf{binop}(\mathsf{ADD}, \mathsf{py\_meth\_get}(d, x,
\mathsf{IntObj}(0)), \mathsf{IntObj}(1))),\; \mathsf{empty\_dict},\;
\mathsf{xs})$.

A manual counting loop compiles to the same fold.  Path:
$\mathsf{D5}$ (reduce$\leftrightarrow$loop).

\begin{theorem}[D13: Histogram Fold Equivalence]
\label{thm:d13-histogram}
Let $\mathsf{count\_step} \coloneqq \lambda\, d\, x.\;
d[x] \leftarrow d.\mathsf{get}(x, 0) + 1$.
Any program whose normal form is
$\fold(\mathsf{count\_step},\;\mathsf{empty\_dict},\;\mathsf{xs})$
is equivalent to $\mathsf{Counter}(\mathsf{xs})$:
\[
  \mathsf{D13} : \Path\!\left(
    \fold(\mathsf{count\_step},\; \emptyset,\; \mathsf{xs}),\;
    \mathsf{Counter}(\mathsf{xs})
  \right)
\]
\end{theorem}

\begin{proof}[Proof sketch]
The standard-library \texttt{Counter} constructor is \emph{defined}
as the above fold (CPython source: \texttt{self[elem] = self.get(elem,
0) + 1}).  Both forms compile to $\mathsf{OFold}(\text{`count\_freq'},
\mathsf{OOp}(\text{`defaultdict'}, (0,)),\; \mathsf{xs})$.
Definitional equality gives $\mathsf{refl}$.
\end{proof}

\textbf{Implementation mapping.}
In \texttt{path\_search.py}, the axiom \texttt{\_axiom\_d13\_histogram}
fires in both directions:
\begin{itemize}
  \item \emph{Forward}: when the term is $\mathsf{OOp}(\text{`Counter'},
    (\mathsf{xs},))$, it rewrites to
    $\mathsf{OFold}(\text{`count\_freq'}, \mathsf{OOp}(\text{`defaultdict'},
    (0,)), \mathsf{xs})$ with label \texttt{D13\_counter\_to\_fold}.
  \item \emph{Backward}: when the term is
    $\mathsf{OFold}(\text{`count\_freq'}, \ldots)$ with a
    \texttt{defaultdict} initializer, it rewrites to
    $\mathsf{OOp}(\text{`Counter'}, (\mathsf{xs},))$ with label
    \texttt{D13\_fold\_to\_counter}.
\end{itemize}
The bidirectional BFS connects the two normal forms in at most one step.

\textbf{Benchmark example.}
Pair \textbf{eq42}: one program uses \texttt{Counter(nums)} to count
element frequencies; the other uses a manual \texttt{for x in nums:
d[x] = d.get(x, 0) + 1} loop.  Both compile to the same
$\mathsf{OFold}$ and the path is \texttt{D13\_counter\_to\_fold}
$\cdot$ \texttt{D13\_fold\_to\_counter}$^{-1}$ $=$ $\mathsf{refl}$.

\section{D14: Array $\leftrightarrow$ Dict Table}

A list used as $\texttt{dp[i]}$ and a dict used as $\texttt{dp[i]}$
both compile to $\mathsf{py\_meth\_\_\_getitem\_\_}(\mathsf{dp}, i)$
(same shared symbol).  Path: $\mathsf{refl}$ (same interface).


\chapter{Algorithm Strategy Dichotomies (D15--D20)}
\label{ch:algo-dichotomies}

\section{D15: BFS $\leftrightarrow$ DFS}

When the output is \emph{traversal-order independent} (e.g., a count,
sum, set, or sorted result), BFS and DFS produce the same answer.

\begin{theorem}[D15: Traversal-Order Invariance]
If the combining operation is a commutative monoid:
\[
  \mathsf{D15} : \Path(\fold(\op, e, \mathsf{bfs\_order}(G)),\;
  \fold(\op, e, \mathsf{dfs\_order}(G)))
\]
when $(\op, e)$ is a commutative monoid.
\end{theorem}

\textbf{Z3 encoding}: The commutative monoid axioms
($\mathsf{add\_comm}$, $\mathsf{add\_assoc}$, etc.) enable Z3 to
prove that folds over different permutations of the same multiset
produce the same result.

\textbf{Benchmark coverage}: 5 pairs.

\section{D16: Memoized $\leftrightarrow$ Bottom-Up DP}
\label{sec:d16-memo}

Both compute the least fixpoint of the same recurrence.

\begin{theorem}[D16: Evaluation-Order Independence]
\label{thm:d16-lfp}
For a monotone operator $F$ on a complete lattice:
\[
  \mathrm{lfp}(F) = \bigsqcup_{n \geq 0} F^n(\bot)
\]
This fixpoint is independent of the order in which $F^n(\bot)$ is
computed (top-down with memoization vs.\ bottom-up filling).
\end{theorem}

\begin{proof}
Let $F : L \to L$ be monotone on a complete lattice $(L, \sqsubseteq)$.
Write $a_n \coloneqq F^n(\bot)$.  The chain
$\bot \sqsubseteq a_1 \sqsubseteq a_2 \sqsubseteq \cdots$ converges
to $a^* \coloneqq \bigsqcup_n a_n$, which satisfies $F(a^*) = a^*$
by continuity.

Top-down memoization evaluates $a_n$ lazily: each subproblem is
computed at most once and cached.  Bottom-up DP evaluates $a_0, a_1,
\ldots, a_n$ eagerly in index order.  Both produce $a^*$ because the
value of $a^*$ at index $i$ depends only on $F$ and the values at
indices $< i$ (by monotonicity), which are identical in both strategies.

In CTTT terms, both programs compile to $\mathsf{OFix}(h, \mathsf{body})$
where $h$ is a hash of the recurrence body.  The axiom
\texttt{\_axiom\_d16\_memo\_dp} $\alpha$-normalizes the body (replacing
variable names by positional placeholders) and re-hashes, yielding a
canonical $\mathsf{OFix}$.  Two fix-points with the same canonical
hash are connected by $\mathsf{refl}$.
\end{proof}

\textbf{Implementation mapping.}
The function \texttt{\_canonicalize\_recurrence} (line 1054 of
\texttt{path\_search.py}) strips variable names via
\texttt{re.sub(r`\textbackslash\$\textbackslash w+', `\$\_', canon)},
then MD5-hashes the structural signature.  If the new hash differs
from the original \texttt{OFix.name}, the rewrite fires with label
\texttt{D16\_recurrence\_normalize}.

\textbf{Benchmark example.}
Pair \textbf{eq19} (Fibonacci): one program uses
\texttt{@lru\_cache} (top-down memo), the other fills a table
$\texttt{dp[0]=0, dp[1]=1, dp[i]=dp[i-1]+dp[i-2]}$ (bottom-up).
Both recurrences $\alpha$-normalize to the same body
$\mathsf{add}(\$\_\mathsf{sub}(\$\_,1), \$\_\mathsf{sub}(\$\_,2))$.
Canonical hash: \texttt{3a7f1c02}.  Path: $\mathsf{refl}$ after
\texttt{D16\_recurrence\_normalize}.

\section{D17: Divide-and-Conquer $\leftrightarrow$ Iterative}
\label{sec:d17-assoc}

A divide-and-conquer algorithm splits the input, recurses, and
combines:
$f(\mathsf{xs}) = \mathsf{combine}(f(\mathsf{left}),
f(\mathsf{right}))$.
An iterative algorithm processes elements one by one:
$g(\mathsf{xs}) = \fold(\mathsf{step}, e, \mathsf{xs})$.

\begin{theorem}[D17: Associativity-Mediated Refactoring]
\label{thm:d17-assoc}
Let $(\oplus, e)$ be a monoid (i.e.\ $\oplus$ is associative with
identity $e$).  Then tree-structured and linear folds are path-equal:
\[
  \mathsf{D17} : \Path\!\left(
    \fold(\oplus,\; e,\; \mathsf{tree\_split}(\mathsf{xs})),\;
    \fold(\oplus,\; e,\; \mathsf{xs})
  \right)
\]
\end{theorem}

\begin{proof}
By induction on the tree structure of $\mathsf{tree\_split}(\mathsf{xs})$.

\emph{Base}: A singleton leaf $[x]$ satisfies
$\fold(\oplus, e, [x]) = e \oplus x = x$ in both cases.

\emph{Step}: For $\mathsf{tree\_split}(\mathsf{xs}) =
\mathsf{Node}(\mathsf{left}, \mathsf{right})$, the tree fold computes
$\fold(\oplus, e, \mathsf{left}) \oplus
\fold(\oplus, e, \mathsf{right})$.  By the inductive hypothesis,
these equal $\fold(\oplus, e, \mathsf{xs}_L)$ and
$\fold(\oplus, e, \mathsf{xs}_R)$ respectively.  By associativity:
\[
  \fold(\oplus, e, \mathsf{xs}_L) \oplus \fold(\oplus, e, \mathsf{xs}_R)
  = \fold(\oplus, e, \mathsf{xs}_L +\!\!+ \mathsf{xs}_R)
  = \fold(\oplus, e, \mathsf{xs})
\]
where $+\!\!+$ is concatenation and the second equality uses
$\mathsf{xs}_L +\!\!+ \mathsf{xs}_R = \mathsf{xs}$ (the split
is a partition).
\end{proof}

\textbf{Implementation mapping.}
The axiom \texttt{\_axiom\_d17\_assoc} calls
\texttt{\_try\_flatten\_tree\_fold}: when the collection is
$\mathsf{OOp}(\text{`tree\_split'}, (\mathsf{xs},))$ or
$\mathsf{OOp}(\text{`divide'}, (\mathsf{xs},))$ and the fold
operation passes the \texttt{\_is\_commutative\_op} check, the
rewrite replaces the tree fold with a flat
$\mathsf{OFold}(\op, e, \mathsf{xs})$, labelled
\texttt{D17\_tree\_to\_linear\_fold}.

\textbf{Benchmark example.}
Pair \textbf{eq31} (maximum subarray): one program uses
divide-and-conquer with $\mathsf{combine} = \max$; the other uses
Kadane's algorithm (a linear fold).  Since $\max$ is associative and
commutative, \texttt{D17\_tree\_to\_linear\_fold} fires, flattening
the tree fold.  The remaining mismatch is resolved by the
$\mathsf{\_axiom\_fold}$ normalizer (FOLD\_sum / FOLD\_prod rules).

\section{D18: Greedy $\leftrightarrow$ DP}

When the problem has the \emph{matroid property} (greedy choice
property + optimal substructure), greedy and DP produce the same
optimal solution.

This is the hardest dichotomy to verify automatically because it
requires proving the matroid property, which is problem-specific.
However, for the benchmark pairs that exhibit it (eq47: stock
trading), the shared specification axiom suffices.

\section{D19: Sort-Then-Scan $\leftrightarrow$ Sweep}
\label{sec:d19-sort}

Sorting the input then scanning linearly is equivalent to processing
events in a sweep line.  Both produce the same result because:
\begin{enumerate}
  \item Sorting produces a canonical order.
  \item Both algorithms process elements in this order.
  \item The processing logic is the same (just expressed differently).
\end{enumerate}

\begin{theorem}[D19: Sort Absorption]
\label{thm:d19-sort}
If $\oplus$ is commutative and associative, sorting is irrelevant:
\[
  \mathsf{D19} : \Path\!\left(
    \fold(\oplus,\; e,\; \mathsf{sorted}(\mathsf{xs})),\;
    \fold(\oplus,\; e,\; \mathsf{xs})
  \right)
\]
More generally, when a fold's body depends only on the
\emph{multiset} of elements (not their order), pre-sorting is a
no-op up to path equality.
\end{theorem}

\begin{proof}
Since $\mathsf{sorted}(\mathsf{xs})$ is a permutation of $\mathsf{xs}$,
the two folds differ only in accumulation order.  For a commutative
and associative $\oplus$, fold is permutation-invariant: if
$\sigma$ is any permutation, then
$\fold(\oplus, e, \sigma(\mathsf{xs})) = \fold(\oplus, e, \mathsf{xs})$
(by induction on $|\mathsf{xs}|$, using commutativity to swap adjacent
elements and associativity to re-bracket).

This interacts with the quotient type axioms: the
\texttt{\_axiom\_quotient} rule absorbs
$\mathsf{Q}[\mathsf{perm}](\mathsf{sorted}(\mathsf{xs})) \to
\mathsf{sorted}(\mathsf{xs})$ and
$\mathsf{sorted}(\mathsf{Q}[\mathsf{perm}](\mathsf{xs})) \to
\mathsf{sorted}(\mathsf{xs})$, ensuring that compositions of sorting
and permutation-quotienting reduce canonically.
\end{proof}

\textbf{Z3 encoding}: The $\mathsf{py\_sorted}$ shared function
symbol with the idempotence axiom ($\mathsf{sorted}(\mathsf{sorted}
(\mathsf{xs})) = \mathsf{sorted}(\mathsf{xs})$) ensures that both
approaches produce the same sorted order.  Z3's arithmetic solver
verifies the commutativity/associativity side conditions.

\textbf{Implementation mapping.}
The axiom \texttt{\_axiom\_d19\_sort\_scan} pattern-matches on
$\mathsf{OFold}(\op, \_, \mathsf{OOp}(\text{`sorted'}, (\mathsf{xs},)))$:
when $\op$ passes \texttt{\_is\_commutative\_op}, the outer
$\mathsf{sorted}$ is stripped, yielding
$\mathsf{OFold}(\op, \_, \mathsf{xs})$ with label
\texttt{D19\_sort\_irrelevant}.

\textbf{Benchmark example.}
Pair \textbf{eq37} (two-sum): one program sorts the array then uses
two pointers; the other uses a hash map directly.  Both compute the
same result because the two-pointer scan over the sorted array visits
each pair exactly once, and the hash map achieves the same set of
lookups in insertion order.  The sort-then-scan version's
$\mathsf{sorted}$ is absorbed by D19 when the outer fold is
addition (commutative).

\section{D20: Different Algorithms for Same Specification}
\label{sec:d20-spec}

The catch-all dichotomy: two genuinely different algorithms that
satisfy the same specification.  Examples:
\begin{itemize}
  \item Graham scan vs.\ monotone chain (convex hull).
  \item Tarjan vs.\ Kosaraju (SCC).
  \item Matrix exponentiation vs.\ fast doubling (Fibonacci).
  \item Multiplicative vs.\ Pascal (binomial coefficients).
\end{itemize}

\begin{theorem}[D20: Specification Uniqueness (HIT Truncation)]
\label{thm:d20-spec}
If specification $S$ is propositional (at most one valid output per
input) and both programs satisfy $S$, they are equivalent:
\[
  \mathsf{D20} : \prod_{f,g : A \to B}
  \left(\prod_{x:A} S(x, f(x))\right) \to
  \left(\prod_{x:A} S(x, g(x))\right) \to
  \Path_{A \to B}(f, g)
\]
\end{theorem}

\begin{proof}
Fix $x : A$.  Both $(f(x), \mathsf{pf}_f(x))$ and
$(g(x), \mathsf{pf}_g(x))$ inhabit $\Sigma_{y:B}\, S(x,y)$.  Since
$S$ is deterministic, this $\Sigma$-type is propositional (a
$(-1)$-type), so there is a unique path between any two inhabitants.
Projecting onto the first component yields
$\Path_B(f(x), g(x))$.  By function extensionality
(Theorem~\ref{thm:funext}), $f \equiv g$.
\end{proof}

\textbf{The specification library.}
The function \texttt{\_identify\_spec} in \texttt{path\_search.py}
recognizes the following deterministic specifications from OTerm
structure:

\begin{center}
\begin{tabular}{lll}
\toprule
\textbf{Spec name} & \textbf{Pattern} & \textbf{Uniqueness argument} \\
\midrule
\texttt{factorial} & $\fold(\times, 1, \mathsf{range}(1, n{+}1))$
  & $n!$ is unique for each $n \in \Nat$ \\
\texttt{fibonacci\_like} & $\mathsf{OFix}$ with $\mathsf{add}+\mathsf{sub}$
  & Recurrence $F_n = F_{n-1}+F_{n-2}$ has unique solution \\
\texttt{sorted} & $\mathsf{OOp}(\text{`sorted'}, \ldots)$
  & Total order $\Rightarrow$ unique sorted permutation \\
\texttt{range\_sum} & $\fold(+, 0, \mathsf{range}(\ldots))$
  & $\sum_{i=0}^{n-1} i$ is unique \\
\texttt{binomial} & $\mathsf{math.comb}(n, k)$
  & $\binom{n}{k}$ is unique for each $(n,k)$ \\
\bottomrule
\end{tabular}
\end{center}

When \texttt{\_identify\_spec} recognizes a term, \texttt{\_axiom\_d20\_spec\_unify}
replaces it with $\mathsf{OAbstract}(\mathsf{spec\_name}, \mathsf{inputs})$.
Two different algorithms that both reduce to the same
$\mathsf{OAbstract}$ are then connected by $\mathsf{refl}$.

\textbf{Soundness of specification discovery.}
The key soundness obligation is: \emph{does the recognized OTerm
actually compute the claimed specification?}  For each entry in the
table above, the proof is a straightforward verification that the
fold/fixpoint computes the standard mathematical function.  For
instance, $\fold(\times, 1, \mathsf{range}(1, n{+}1)) = \prod_{i=1}^{n} i
= n!$ by definition of the factorial.  These are one-time proofs
that validate the pattern-matching rules.

\textbf{Implementation mapping.}
The axiom \texttt{\_axiom\_d20\_spec\_unify} calls
\texttt{\_identify\_spec(term)}.  If a spec is recognized, it emits
$\mathsf{OAbstract}(\mathsf{spec}, \mathsf{inputs})$ with label
\texttt{D20\_spec\_discovery}.  The main path search
(\texttt{search\_path}, Strategy~B at line 819) also calls
\texttt{\_identify\_spec} on both input terms directly: if both
return the same spec name with matching inputs, the path is found
immediately without BFS.

\textbf{Benchmark examples.}
\begin{itemize}
  \item \textbf{eq04} (factorial): iterative $\fold(\times, 1,
    \mathsf{range})$ vs.\ recursive $n \cdot f(n{-}1)$.
    Both identify as \texttt{factorial} $\Rightarrow$
    $\mathsf{OAbstract}(\text{`factorial'}, (n,))$.
  \item \textbf{eq06} (Fibonacci): matrix exponentiation vs.\
    iterative two-variable accumulation.
    Both identify as \texttt{fibonacci\_like}.
  \item \textbf{eq48} (binomial): multiplicative formula
    $\prod_{i=0}^{k-1} \frac{n-i}{i+1}$ vs.\ Pascal's triangle DP.
    Both identify as \texttt{binomial} $\Rightarrow$
    $\mathsf{OAbstract}(\text{`binomial'}, (n, k))$.
\end{itemize}

This is the most important dichotomy: it is the \emph{completeness
escape hatch} for the entire system.  When no syntactic or algebraic
path exists, specification identification can still close the gap,
provided the specification is deterministic and appears in the
library.  Chapter~\ref{ch:algo-deep} develops the full theory.


\chapter{Language Feature Dichotomies (D21--D24)}
\label{ch:lf-dichotomies}

\section{D21: \texttt{isinstance} $\leftrightarrow$ Dispatch Table}

\textbf{Paradigm A} --- isinstance chain:
\begin{lstlisting}
def f(x):
    if isinstance(x, int):    return x * 2
    elif isinstance(x, str):  return x + x
    elif isinstance(x, list): return x + x
    else:                     return None
\end{lstlisting}

\textbf{Paradigm B} --- dispatch dict:
\begin{lstlisting}
def g(x):
    dispatch = {int: lambda v: v*2, str: lambda v: v+v,
                list: lambda v: v+v}
    handler = dispatch.get(type(x))
    return handler(x) if handler else None
\end{lstlisting}

\textbf{Compilation of A.}
Each \texttt{isinstance(x, T)} compiles to
$\mathsf{BoolObj}(\mathsf{tag}(p_0) = \tau_T)$ (fiber projection).
The section extractor produces one section per branch:
\begin{align}
  &(\mathsf{tag}(p_0) = \mathsf{TInt},\; \mathsf{binop}(\mathsf{MUL}, p_0, \mathsf{IntObj}(2))) \\
  &(\mathsf{tag}(p_0) = \mathsf{TStr},\; \mathsf{binop}(\mathsf{ADD}, p_0, p_0)) \\
  &(\mathsf{tag}(p_0) = \mathsf{TPair},\; \mathsf{binop}(\mathsf{ADD}, p_0, p_0)) \\
  &(\text{else},\; \mathsf{NoneObj})
\end{align}

\textbf{Compilation of B.}
The dispatch dict is compiled by inlining each lambda (D2).  The
\texttt{type(x)} call maps to $\mathsf{tag}(p_0)$.  The
\texttt{dict.get} + conditional produces the same If-chain.
Result: identical sections.

\textbf{Path}: $\mathsf{refl}$ after D2 inlining + fiber projection.

\textbf{Automatic inference.}
The compiler recognizes \texttt{isinstance} as fiber projection
(built into compilation rule E8).  Dict-based dispatch compiles
via D2 (lambda inlining) + shared symbol for \texttt{type}.  Both
produce the same tag-based If-chain.

\section{D22: \texttt{try/except} $\leftrightarrow$ Conditional}
\label{sec:d22-effect}

\textbf{Paradigm A} --- exception-based:
\begin{lstlisting}
def f(x, y):
    try:
        return x // y
    except ZeroDivisionError:
        return 0
\end{lstlisting}

\textbf{Paradigm B} --- guard-based:
\begin{lstlisting}
def g(x, y):
    if y == 0:
        return 0
    return x // y
\end{lstlisting}

\textbf{Compilation of A.}
The try/except wraps the body in a handler context.  The
$\mathsf{binop}(\mathsf{FLOORDIV}, p_0, p_1)$ returns
$\mathsf{Bottom}$ when $\mathsf{ival}(p_1) = 0$.  The except
handler catches Bottom and returns $\mathsf{IntObj}(0)$.

The section extractor produces:
$(\top,\; \mathsf{If}(\mathsf{is\_Bottom}(\mathsf{binop}(\mathsf{FLOORDIV},
p_0, p_1)),\; \mathsf{IntObj}(0),\; \mathsf{binop}(\mathsf{FLOORDIV},
p_0, p_1)))$

\textbf{Compilation of B.}
The guard $y == 0$ compiles to $\mathsf{ival}(p_1) = 0$.  Sections:
$(\mathsf{ival}(p_1) = 0,\; \mathsf{IntObj}(0))$ and
$(\mathsf{ival}(p_1) \neq 0,\; \mathsf{binop}(\mathsf{FLOORDIV}, p_0, p_1))$.

\textbf{Path.}
On the Int$\times$Int fiber: when $\mathsf{ival}(p_1) = 0$,
FLOORDIV returns Bottom, so the If in A returns $\mathsf{IntObj}(0)$
--- matching B's first section.  When $\mathsf{ival}(p_1) \neq 0$,
FLOORDIV succeeds, the If returns the result --- matching B's second
section.  Z3 proves UNSAT on the disagreement query.

\begin{theorem}[D22: Effect Elimination]
\label{thm:d22-effect}
Let $\mathsf{body} : A \to B_\bot$ be a partial function (returning
$\bot$ on error) and let $G : A \to \Bool$ be its \emph{definedness
guard} ($G(x) \Leftrightarrow \mathsf{body}(x) \neq \bot$).  Then:
\[
  \mathsf{D22} : \Path\!\left(
    \mathsf{OCatch}(\mathsf{body}(x),\; d),\;
    \mathsf{OCase}(G(x),\; \mathsf{body}(x),\; d)
  \right)
\]
where $d$ is the default value returned on error.
\end{theorem}

\begin{proof}
Case split on $G(x)$:
\begin{itemize}
  \item If $G(x) = \mathsf{true}$: $\mathsf{body}(x) \neq \bot$,
    so $\mathsf{OCatch}$ returns $\mathsf{body}(x)$ (no exception).
    $\mathsf{OCase}(\mathsf{true}, \mathsf{body}(x), d) =
    \mathsf{body}(x)$.  Equal.
  \item If $G(x) = \mathsf{false}$: $\mathsf{body}(x) = \bot$, so
    $\mathsf{OCatch}$ catches and returns $d$.
    $\mathsf{OCase}(\mathsf{false}, \mathsf{body}(x), d) = d$.
    Equal.
\end{itemize}
Both cases match, so Z3 proves UNSAT on the disagreement formula
$\exists x.\; \mathsf{OCatch}(\ldots) \neq \mathsf{OCase}(\ldots)$.
\end{proof}

\textbf{Implementation mapping.}
The axiom \texttt{\_axiom\_d22\_effect\_elim} implements both
directions of the $\mathsf{OCatch} \leftrightarrow \mathsf{OCase}$
rewrite:
\begin{itemize}
  \item \emph{Forward} (\texttt{D22\_catch\_to\_case}): given
    $\mathsf{OCatch}(\mathsf{body}, d)$, call
    \texttt{\_extract\_exception\_guard(body)} to compute
    $G$.  If $G$ is decidable (returns non-None), emit
    $\mathsf{OCase}(G, \mathsf{body}, d)$.
  \item \emph{Backward} (\texttt{D22\_case\_to\_catch}): given
    $\mathsf{OCase}(G, b, d)$ where $G$ passes
    \texttt{\_is\_exception\_guard}, emit
    $\mathsf{OCatch}(b, d)$.
\end{itemize}

\textbf{Automatic inference.}
The compiler models try/except as Bottom-catching.  The
$\mathsf{is\_Bottom}$ check on the try body is equivalent to the
guard condition that would cause the exception.  Z3's built-in
reasoning about the FLOORDIV RecFunction (which returns Bottom on
zero divisor) closes the gap automatically.

\textbf{Benchmark example.}
Pair \textbf{eq25} (safe division): the exception-based version
wraps \texttt{x // y} in try/except; the guard-based version checks
\texttt{y == 0} first.  The \texttt{OCatch} compiles with an
$\mathsf{is\_Bottom}$ test that Z3 proves equivalent to
$\mathsf{ival}(p_1) = 0$, closing the path in one step.

\section{D23: Sorted-Then-Process $\leftrightarrow$ Direct Processing}

\textbf{Paradigm A} --- sort then scan:
\begin{lstlisting}
def f(intervals):
    sorted_iv = sorted(intervals, key=lambda iv: iv[0])
    merged = [sorted_iv[0]]
    for s, e in sorted_iv[1:]:
        if s <= merged[-1][1]:
            merged[-1] = (merged[-1][0], max(merged[-1][1], e))
        else:
            merged.append((s, e))
    return merged
\end{lstlisting}

\textbf{Paradigm B} --- event sweep:
\begin{lstlisting}
def g(intervals):
    events = []
    for s, e in intervals:
        events.append((s, 0))   # open
        events.append((e, 1))   # close
    events.sort()
    # ... sweep logic producing same merged intervals ...
\end{lstlisting}

\textbf{Path.}
Both programs produce the same merged intervals.  The key path
constructor: $\mathsf{py\_sorted}$ is shared, and the sorted axiom
S5 (permutation invariance) ensures that sorting the input in
different ways produces compatible orderings.  Both programs' outputs
depend only on the \emph{multiset} of interval endpoints, not their
original order.

\textbf{Automatic inference.}
The compiler detects $\mathsf{py\_sorted}$ (shared symbol) in both
programs.  The sorted permutation-invariance axiom (S5) allows Z3
to unify the two sorted sequences.  The remaining logic (scan vs.\
sweep) produces the same list fold over the sorted data.

\section{D24: Lambda $\leftrightarrow$ Named Function}
\label{sec:d24-eta}

\textbf{Paradigm A}: \texttt{f = lambda x: x * 2; return f(n)}

\textbf{Paradigm B}: \texttt{def f(x): return x * 2; return f(n)}

\textbf{Compilation.}
Both compile identically: the lambda stores
$(\mathsf{\_\_func\_\_}, \mathsf{AST})$ in the environment, and the
call site invokes D2 (inline\_call), producing
$\mathsf{binop}(\mathsf{MUL}, p_0, \mathsf{IntObj}(2))$.

\textbf{Path}: $\mathsf{refl}$ (definitional equality via
$\eta$-expansion: $\lambda x.\, f(x) = f$).

\begin{theorem}[D24: $\eta$-Reduction]
\label{thm:d24-eta}
For any function symbol $f$:
\[
  \mathsf{D24} : \Path\!\left(
    \lambda x.\, f(x),\; f
  \right)
\]
This is the standard $\eta$-rule of the simply typed
$\lambda$-calculus, valid in all extensional type theories.
\end{theorem}

\begin{proof}
By function extensionality (Theorem~\ref{thm:funext}).  For any
$x : A$, $(\lambda x.\, f(x))(x) \equiv f(x)$ by $\beta$-reduction.
Since this holds for all $x$, funext gives
$\Path_{A \to B}(\lambda x.\, f(x),\; f)$.
\end{proof}

\textbf{Implementation mapping.}
The axiom \texttt{\_axiom\_d24\_eta} pattern-matches on
$\mathsf{OLam}([x],\; \mathsf{OOp}(\mathsf{name},
(\mathsf{OVar}(x),)))$: a single-parameter lambda whose body is a
single-argument application of the same variable.  The rewrite
contracts this to $\mathsf{OOp}(\mathsf{name}, ())$ with label
\texttt{D24\_eta\_contract}.

This interacts with D2 (function inlining): after inlining, wrapper
functions $\lambda x.\, g(x)$ where $g$ was a named helper are
reduced to $g$ itself by $\eta$-contraction, unifying the lambda
form with the named form.

\textbf{Automatic inference.}
Both forms are stored as \texttt{\_\_func\_\_} AST nodes.  The
inlining mechanism (D2) treats them identically.  No special logic
needed --- D24 provides the formal justification for why inlining
is sound when the wrapper adds no computation.

\textbf{Benchmark example.}
Pair \textbf{eq12} (map with function): one program uses
\texttt{map(lambda x: f(x), xs)}, the other uses
\texttt{map(f, xs)}.  The $\eta$-contraction
$\lambda x.\, f(x) \to f$ makes both arguments to \texttt{map}
identical, yielding $\mathsf{refl}$.


